\begin{table}[htbp]
\centering
\caption{The Lorentzian linewidth is measured, not fitted to identification. For each acquisition a single width shared by all detected bands was fitted simultaneously with the band amplitudes, band centres being free to move by at most one pixel.}
\label{tab:si-gamma}
\small
\begin{tabular}{lrrrrr}
\toprule
Agent & Acquisitions & $\Gamma$ median (\wn) & Range (\wn) & $R^2$ median & Bands fitted \\
\midrule
GD & 4 & 18.9 & 18.4--19.8 & 0.92 & 19 \\
GF & 3 & 17.4 & 17.4--17.7 & 0.91 & 20 \\
VX & 1 & 38.9 & 38.9--38.9 & 0.89 & 17 \\
HD & 1 & 48.4 & 48.4--48.4 & 0.81 & 7 \\
HN1 & 3 & 21.7 & 21.6--21.7 & 0.91 & 16 \\
HN2 & 3 & 9.3 & 9.1--9.3 & 0.98 & 7 \\
HN3 & 3 & 22.5 & 22.3--22.7 & 0.87 & 6 \\
A-230 & 1 & 17.8 & 17.8--17.8 & 0.82 & 10 \\
A-232 & 3 & 24.4 & 24.2--25.5 & 0.92 & 18 \\
A-234 & 4 & 29.8 & 21.6--30.0 & 0.91 & 17 \\
A-242 & 2 & 33.7 & 32.2--35.2 & 0.91 & 18 \\
\bottomrule
\end{tabular}

\begin{flushleft}\footnotesize Of the 32 genuine acquisitions, 28 yielded a converged fit ($R^2 > 0.8$); the remaining 4 did not and are excluded from this table. Over the 28 the median is 21.7\,\wn{} with an interquartile range of 18.2--26.5\,\wn{}. We adopt $\Gamma$ = 22\,\wn{}. The extremes are not trustworthy: HN2 offers a single band, and HD and VX a single acquisition each. This is an apparent width that includes the instrument response function. No identification result was consulted in obtaining it.\end{flushleft}
\end{table}
